% chapters/15/rangkuman.tex
\bagian{Rangkuman}

\begin{rangkumanbox}
\begin{itemize}
  \item Bunyi menjangkau pemain tanpa menuntut perhatian matanya. Karena itu,
        setiap keterangan yang disampaikan melalui bunyi sebaiknya disampaikan
        pula secara terlihat, agar permainan tetap dapat dipahami oleh pemain
        yang pendengarannya terganggu.
  \item Ada dua jenis bunyi yang perlu dibedakan sejak awal. Musik latar
        berjalan terus dan tidak boleh bertumpuk; efek suara dipanggil menurut
        kejadian dan justru perlu dapat bertumpuk.
  \item Satu \code{AudioSource} tidak boleh dipakai untuk keduanya, sebab setiap
        efek suara akan menghentikan musik yang sedang berjalan.
  \item \code{AudioClip} adalah datanya, \code{AudioSource} adalah pemutarnya.
        Letak bunyi di dalam ruang ditentukan oleh \code{AudioSource}, bukan
        oleh berkasnya.
  \item \istilah{Spatial Blend} bernilai nol untuk musik latar dan antarmuka,
        dan bernilai satu untuk bunyi benda di dalam level.
  \item \code{PlayOneShot} dipakai untuk efek suara, sebab beberapa bunyi dapat
        berbunyi bersamaan. \code{Play} dipakai untuk musik latar, sebab hanya
        satu musik yang boleh berjalan.
  \item Bunyi dipusatkan pada satu skrip, agar hanya ada satu musik yang
        berjalan dan perbandingan kekuatannya tetap sama di mana pun bunyi itu
        dipanggil.
  \item Bunyi dipanggil pada peristiwa, bukan pada keadaan. Memanggilnya
        berdasarkan keadaan menghasilkan bunyi yang berulang setiap bingkai.
  \item Bunyi dipanggil di \code{Update}, bukan di \code{FixedUpdate}, sebab
        penundaan dua sampai tiga bingkai sudah terasa pada bunyi yang pendek.
  \item Kekuatan bunyi pada \istilah{AudioMixer} dinyatakan dalam desibel.
        Menurunkan 6 dB berarti mengurangi kekuatannya menjadi separuh, tetapi
        perubahannya tidak terdengar sebesar itu.
  \item Selisih sekitar 10 dB antara efek suara dan musik latar adalah titik
        awal yang baik. Menyeimbangkan bunyi lebih sering berarti menurunkan
        yang terlalu kuat daripada menaikkan yang terlalu pelan.
  \item Menggeser nada sedikit setiap kali mencegah bunyi yang berulang
        terdengar seperti mesin.
  \item Partikel adalah pasangan yang terlihat bagi bunyi, dan dipanggil pada
        peristiwa yang sama.
  \item Hak pakai berkas bunyi lebih berat akibatnya daripada hak pakai gambar,
        sebab musik lebih mudah dikenali. Catat sumber dan lisensinya sejak
        mengunduh.
  \item Bab ini menuntaskan \SUBCPMK{4.2}. Seluruh bagian yang terlihat dan
        terdengar dari level Anda sudah lengkap; yang tersisa adalah mengujinya
        pada Bab~XVI.
\end{itemize}
\end{rangkumanbox}
